Antichains of heaps appear frequently in this paper. We highlight one property of maximal antichains now. Recall
that an \emph{ideal} in a poset $P$ is a set $\cali$ such that if $y\in \cali$ and $x<y$ for $x,y\in P$ then $x\in
\cali$ ; dually, a \emph{filter} in $P$ is a set $\calf$ such that if $y\in \calf$ and $x>y$ for $x,y\in P$ then
$x\in \calf$. Any set $A\se P$ \emph{generates} an ideal 
\[ \cali_A:=\{i\in H:i\le j\text{ for some $j\in A$}\} \]
and a filter 
\[ \calf_A:=\{i\in H:i\ge j\text{ for some $j\in A$}\}.\] 
If $A$ is an antichain, then any element in the set $P\setminus (\cali_A\cup\calf_A)$, if it exists, can be
adjoined to $A$ to form a larger antichain. This implies the following:

\begin{lemma}
    \label{lemm:maxAntichain} If $A$ is a maximal antichain in $P$, then $P=\cali_A\cup \calf_A$.
\end{lemma}
